% Introduction (v9, 2026-08-19). natbib + paper/references.bib. No em-dashes, no numbers.
% v9 = v8 trimmed ~35%: every paragraph tightened, redundant clauses merged.

\section{Introduction}
\label{sec:intro}

Vision-language-action models (VLAs) have made large-scale pretraining the dominant recipe for robot
manipulation~\citep{rt2, openvla, pi0}. The pretrained backbone provides priors strong enough that a
new task rarely starts from scratch; in practice, a VLA is deployed by fine-tuning it on
demonstrations of the target task.

The limitation of this recipe lies in the data it imitates. Teleoperated demonstrations interleave
expert behavior with hesitation, detours, and failure, and their quality varies even within a single
trajectory~\citep{grrl, pistar06}. An imitation objective treats every transition as an equally valid
target, so collecting more data reproduces the same average; filtering discards the scarce experience
the recipe depends on, and online refinement is costly and unsafe at this scale.

Can an optimal policy be extracted from suboptimal human data with a VLA-scale model?
Temporal-difference (TD) learning promises precisely this, stitching high-value segments of mediocre
trajectories into behavior no single demonstration exhibits~\citep{cql, iql}. But a VLA is not the
policy this machinery was designed for: it predicts a fixed, long action chunk~\citep{act, pi0}, it
faces long horizons where one-step TD struggles~\citep{qc}, and its backbone cannot be trained
jointly with a critic at practical cost.

Recent attempts stop short of this promise. Most reshape imitation with outcome
signals~\citep{pistar06, arfm, arm, reinbot, nora15}, inheriting both the stability and the ceiling of
supervised learning. The few that learn a value keep it away from the policy: it filters
data~\citep{grrl}, rescores the policy's own samples~\citep{vgps, deas, robomonkey}, or scores
commitments fixed in advance~\citep{corft}. These are structural dead ends. Rescoring cannot exceed
its proposals and its argmax harvests critic error~\citep{emaq}; a fixed-chunk critic cannot tell
skill from luck in closed-loop demonstrations and overvalues commitments that merely happened to
work~\citep{dqc}. What matters is not only how well the value is learned, but \emph{where it is
allowed to act}.

Skilled motor behavior shows where. A human hand crosses free space ballistically and tightens into
feedback control only when contact matters~\citep{woodworth1899, todorov2002}; even corrections
arrive as short open-loop segments with intermittent replanning~\citep{gawthrop2011}. Acting well is
the judgment of how long to commit, and a fixed chunk denies a VLA exactly that. We let the value
supply it through the two channels that stay inside the data: deciding, state by state, how much of
each predicted chunk to execute, and retraining the policy toward what the value prefers. Concretely,
a prefix-valued transformer critic~\citep{acsac, aqc} over frozen VLA representations, trained with
conservative in-sample targets~\citep{iql, deas}, drives state-dependent commitment and policy
extraction. The result surpasses supervised fine-tuning from the very same demonstrations, without a
single additional rollout.

\paragraph{Contributions.}
(i)~A structural analysis of why filtering, test-time rescoring, and fixed commitments cannot surpass
the policy they start from, each failure mode confirmed in controlled
experiments~\citep{emaq, dqc, park2024bottleneck}. (ii)~The principle that follows and its
instantiation: a value acting through in-data channels, state-dependent commitment and policy
extraction, as a complete offline pipeline on frozen VLA representations. (iii)~A component-isolated
stability recipe for chunk-space TD at VLA scale. (iv)~End-to-end validation on offline RL
benchmarks~\citep{ogbench} and VLA-scale manipulation~\citep{robocasa}, using only the demonstrations
supervised fine-tuning already uses.
